% Amalgamation along units for orthogonal monoids and the induced wdim formula.

\begin{theorem}[正交幺半群沿单位子幺半群并合的 Stokes 维数公式]\label{thm:cdim-orthogonal-amalgamated-wdim}
设 \(M_1,M_2\) 为正交幺半群，并存在同一可逆子幺半群 \(U\) 及单态
\(
\iota_i:U\hookrightarrow U(M_i)\subseteq M_i
\)
（\(i=1,2\)）。
令 \(M\) 为交换幺半群范畴中由 \(\iota_1,\iota_2\) 诱导的并合（推送出）：
\[
M:=M_1\amalg_{U}M_2.
\]
记 \(U_{\mathrm{free}}\) 为 \(U\) 的自由阿贝尔部分，则 \(M\) 仍为正交幺半群，并且满足
\[
\boxed{\;
\mathrm{wdim}(M)=\mathrm{wdim}(M_1)+\mathrm{wdim}(M_2)-\rank_{\ZZ}U_{\mathrm{free}}.\;}
\]
\end{theorem}
\begin{proof}
由正交假设，存在分解
\[
M_i\cong \ZZ^{u_i}\oplus \NN^{v_i}\oplus F_i
\qquad(i=1,2),
\]
其中 \(F_i\) 为有限交换群。
因此其单位群满足
\(
U(M_i)\cong \ZZ^{u_i}\oplus F_i
\)，并且群完成满足
\(
G(M_i)\cong \ZZ^{u_i+v_i}\oplus F_i
\)。
\par
由于 \(\iota_i(U)\subseteq U(M_i)\)，并合仅在单位方向施加识别关系；而 \(\NN^{v_1}\) 与 \(\NN^{v_2}\) 的正锥部分在并合中保持直和（否则将迫使非单位元素在并合中可逆，与正交结构矛盾）。
因此 \(G(M)\) 的无挠秩满足
\[
\rank_{\ZZ}G(M)_{\mathrm{free}}=(u_1+v_1)+(u_2+v_2)-\rank_{\ZZ}U_{\mathrm{free}},
\]
而单位群的无挠秩满足
\[
\rank_{\ZZ}U(M)_{\mathrm{free}}=u_1+u_2-\rank_{\ZZ}U_{\mathrm{free}}.
\]
从而
\[
\rank_{\ZZ}\bigl(G(M)/U(M)\bigr)=v_1+v_2.
\]
由定理 \ref{thm:cdim-stokes-character-contraction-general-monoid} 的恒等式
\(
\mathrm{wdim}(N)=\rank_{\ZZ}U(N)+\frac12\rank_{\ZZ}(G(N)/U(N))
\)
对 \(N=M,M_1,M_2\) 分别应用并代入上述秩计算，即得
\[
\mathrm{wdim}(M)
=\Bigl(u_1+u_2-\rank U_{\mathrm{free}}\Bigr)+\frac12(v_1+v_2)
=\mathrm{wdim}(M_1)+\mathrm{wdim}(M_2)-\rank U_{\mathrm{free}}.
\]
证毕。
\end{proof}

\endinput

